\section{Mathematical Foundations}
\label{sec:math}
% ====================================================================

\subsection{Problem Formulation}

We model the parkour task as a Partially Observable Markov Decision
Process (POMDP) $(\mathcal{S}, \mathcal{A}, \mathcal{O}, T, R, \gamma)$
with state space $\mathcal{S}$, action space $\mathcal{A}$,
observation space $\mathcal{O}$, transition dynamics
$T(s'\mid s,a)$, reward $R(s,a)$, and discount factor $\gamma \in [0,1)$.

The agent receives partial observation $o_t \in \mathcal{O}$
and selects target joint position offsets
$a_t \in \mathcal{A} = \R^{12}$.
Observations decompose into:
\begin{itemize}
	\item \textbf{Proprioception} $o_\text{prop} \in \R^{53}$: base
	      angular velocity, projected gravity, velocity and heading
	      commands, joint positions, joint velocities, previous
	      actions, and foot contacts.
	\item \textbf{Exteroception} $o_\text{ext}$: scandot
	      heightmaps ($\R^{132}$) for the teacher, or depth images
	      for the student.  Scandots~\citep{agarwal2023legged} sample terrain
	      height on a $12 \times 11$ grid in the robot's body frame,
	      recording the height difference between the base and the
	      surface, clipped to $[-1, 1]$.
\end{itemize}

The objective is to find a policy $\pi$ maximizing the expected
discounted return
$J(\pi) = \E_\pi\!\left[\sum_{t=0}^{\infty} \gamma^t R(s_t, a_t)\right]$.
In a POMDP, the optimal policy generally depends on the full belief
state.  In practice, we approximate this with a recurrent architecture
(Section~\ref{sec:student}) that integrates observations over time.

\subsection{PPO}
\label{sec:ppo}

The teacher policy $\pi_\theta$ is trained with Proximal Policy
Optimization (PPO)~\citep{schulman2017proximal}.
Policy gradient methods optimize $J(\pi_\theta)$ by estimating
$\nabla_\theta J$ from trajectory samples, but naively updating
$\theta$ with large steps can collapse performance.
Trust Region Policy Optimization (TRPO)~\citep{schulman2015trust}
constrains the KL divergence between successive policies; PPO
approximates this constraint with a clipped surrogate:
\begin{equation}
	L^{\text{CLIP}} = \E_t\!\left[\min\!\left(
		r_t \hat{A}_t,\;
		\clip(r_t, 1{-}\epsilon, 1{+}\epsilon)\,\hat{A}_t
		\right)\right],
	\label{eq:ppo}
\end{equation}
where $r_t = \pi_\theta(a_t|o_t) / \pi_{\theta_{\text{old}}}(a_t|o_t)$
is the probability ratio.
The advantage estimate $\hat{A}_t$ controls the effect of the gradient.  GAE~\citep{schulman2016high} defines
\begin{equation}
	\hat{A}_t^{\text{GAE}(\gamma,\lambda)}
	= \sum_{l=0}^{\infty} (\gamma\lambda)^l \,\delta_{t+l},
	\quad
	\delta_t = r_t + \gamma V(s_{t+1}) - V(s_t),
	\label{eq:gae}
\end{equation}
where $\delta_t$ is the one-step residual and $\lambda\in[0,1]$.

When $\hat{A}_t > 0$ the clip prevents $r_t$ from exceeding
$1+\epsilon$; when $\hat{A}_t < 0$ it prevents $r_t$ from dropping
below $1-\epsilon$.


The teacher's reward combines task progress with locomotion
regularization. Table~\ref{tab:reward} lists all terms and scales.

\subsection{DAgger Distillation}
\label{sec:dagger}

The student is trained via Dataset Aggregation
(DAgger)~\citep{ross2011reduction}.  Behavioral cloning on a fixed
expert dataset suffers error compounding, but DAgger reduces this by iteratively
collecting data under the student's own distribution while labeling
with the expert.  The student observes $o_\text{prop}$ and a depth
image and predicts $a_\text{student}$, but the teacher's action
$a_\text{teacher}$ is executed.  The loss is:
\begin{equation}
	\mathcal{L} =
	\underbrace{\|a_\text{teacher} - a_\text{student}\|_2}_{\text{action loss}}
	+ \underbrace{\|y_\text{oracle} - \hat{y}\|_2}_{\text{yaw loss}},
	\label{eq:dagger_loss}
\end{equation}
where $\hat{y} \in \R^2$ is the student's predicted heading and
$y_\text{oracle}$ is ground truth heading.

\subsection{Monocular Depth Estimation}
\label{sec:depth_math}

Recovering depth from a single RGB image is an ill-posed problem:
infinitely many 3D scenes can produce the same 2D projection.
Monocular depth estimators resolve this ambiguity by learning
geometric priors from large-scale data.  Given an image
$I \in \R^{H \times W \times 3}$, the model predicts a depth map
$\hat{d} = f_\phi(I) \in \R^{H \times W}$.

\subsection{Depth Signal Quality}
\label{sec:dsq_math}

To measure how well the estimated depth preserves terrain geometry,
we define Depth Signal Quality (DSQ) as the Pearson correlation
between ground-truth and estimated depth, averaged over environments
and time:
\begin{equation}
	\mathrm{DSQ} = \frac{1}{T} \sum_{t=1}^{T}
	\frac{1}{N} \sum_{i=1}^{N}
	\rho(d^{\mathrm{GT}}_i,\, d^{\mathrm{DA2}}_i),
	\label{eq:dsq}
\end{equation}
where $\rho(\cdot,\cdot)$ denotes the Pearson correlation computed
over all pixels $p$ within each frame.

Pearson correlation is invariant to affine transformations of either
argument, i.e. $\rho(\alpha x + \beta,\, y) = \rho(x, y)$ for any
$\alpha > 0$.  This makes DSQ insensitive to the scale--shift
ambiguity inherent in monocular depth.
